\section{Related Work}

\noindent\textbf{Memory for Self-Evolving Agents.} 
Learning from past experiences as procedural memory~\citep{wu2025human, wei2025evo, shen2026decocted, hu2025memory, huang2026rethinking, zhang2024working} is a central mechanism for developing self-evolving agents~\citep{gao2025survey, fang2025comprehensive}. 
The central challenge is to encode interaction histories into reusable and retrievable representations.
Case-based representations are the most concrete form in this research line: they store experiences in minimally processed formats, allowing past histories to be replayed directly or reused as in-context exemplars, such as raw trajectories~\citep{zheng2023synapse, DBLP:journals/corr/abs-2508-16153, wu2025comemagent} and abstracted query–response pairs~\citep{zhao2024expel, islam-etal-2024-mapcoder}. Another line of work abstracts experiences into higher-level knowledge that is editable, auditable, and composable, reducing reliance on long trajectory replay and improving both cross-task generalization and efficiency. Such strategy-based memory typically consists of reusable workflows~\citep{wang2025agent, DBLP:journals/corr/abs-2507-06229}, distilled insights~\citep{ouyang2026reasoningbank, huang-etal-2025-r2d2, DBLP:journals/corr/abs-2509-04439}, and recurring patterns~\citep{yang2024buffer, kim-etal-2025-principles}.
Recently, skills~\citep{wang2025inducing, kuroki2025agent, DBLP:journals/corr/abs-2602-08004, DBLP:journals/corr/abs-2602-12670, DBLP:journals/corr/abs-2602-02474, yang2026autoskillexperiencedrivenlifelonglearning, alzubi2026evoskill, liang2026skillnet} have emerged as a new agent-native form of memory and an orchestrable capability layer, owing to their modularity and ease of customization. Anthropic conceptualizes each skill as a folder containing instructions, scripts, and supporting resources~\citep{anthropic_agent_skills_overview}, which has become the most widely adopted design in the current community. Our work follows this design philosophy, simplifying the setting for research purposes by representing each skill as a single \textit{Markdown} file.

\noindent\textbf{Learning Memory and Skill Curation with RL.} 
Training LLM-based agent systems with memory capabilities using RL has become a growing research direction. One research line targets training for long-context management with predefined operations such as compaction~\citep{zhou2026mem, yu2026memagent, wang2025mem}. Another interesting area focuses more on memory utilization and management by learning additional memory tool-calls~\citep{DBLP:journals/corr/abs-2508-19828, DBLP:journals/corr/abs-2508-16629, DBLP:journals/corr/abs-2510-12635} or training policies for different stages, such as memory retrieval~\citep{zhang2026memrl}. More recently, RL has been applied at various stages of agent skill development. Specifically, SkillRL~\citep{xia2026skillrl} and D2Skill~\citep{tu2026dynamic} teach smaller models to use skills curated from powerful LLMs in an iterative manner. ARISE~\citep{Li2026ARISEAR} trains a shared policy operating both as skill retriever and worker, with heuristics for skill management. 
Recent studies have begun to train agents for memory or skill curation~\citep{DBLP:journals/corr/abs-2512-17102, DBLP:journals/corr/abs-2602-10652}, but their supervision is mostly restricted to local adaptation within short task streams. This favors immediately useful operations such as skill insertion, while offering limited signal for complex management operations, such as revising outdated skills and deleting harmful ones. \ours{} instead formulates skill curation as a long-horizon, executor-grounded learning problem. We group related tasks into training instances and combine downstream task outcomes with intermediate rewards, turning delayed and indirect feedback into learning signals for skill curation.